\chapter{Discussion}
\label{chap:final-discussion}

\section{Summary of the Main Findings}

The overarching objective of this dissertation was to develop, evaluate, and rigorously compare a \acrshortpl{cnn}-based approach for metal artifact reduction in \acrshort{ct}, prioritizing anatomical preservation, robustness to domain variation, and interpretability of the correction. The internal evaluation results support the primary hypothesis of structural improvement in the Artifact Zone relative to the corrupted input: across the six runs (two architectures, three seeds), \acrshort{ssim} in this zone increased by $0{.}322$ to $0{.}328$ (\autoref{tab:roi-delta-ssim}). For the run selected for external evaluation (Single Model, seed 2024), this improvement was statistically significant at the patient level ($\Delta$\acrshort{ssim}$=0{.}3279$, 95\% BCa CI $[0{.}2963;0{.}3548]$, $p_{\mathrm{Holm}}=2{.}73\times10^{-12}$; \autoref{tab:roi-artifact-improvement-ssim}), with an equivalent pattern for the Dual Model ($\Delta$\acrshort{ssim}$=0{.}3266$, $p_{\mathrm{Holm}}=2{.}73\times10^{-12}$). This localized structural improvement was accompanied by a consistent and equally significant global improvement: across all six runs, reconstruction increased global \acrshort{psnr} from $27{.}39$~dB to $34{.}95$--$35{.}20$~dB, global \acrshort{ssim} from $0{.}8361$ to $0{.}9421$--$0{.}9442$, and reduced global \acrshort{mae} from $36{.}6$~HU to $17{.}3$--$17{.}8$~HU (\autoref{tab:global-ssim}, \autoref{tab:global-psnr}, \autoref{tab:global-mae}). At the patient level, the secondary global-improvement tests were significant for all three metrics in every run ($p<0{.}001$; \autoref{tab:delta-ssim-paciente}, \autoref{tab:delta-psnr-paciente}, \autoref{tab:delta-mae-paciente}), with 95\% confidence intervals excluding the null value.

Descriptive slice-level inspection (\autoref{fig:global-improvement-histograms}) showed that unfavorable cases were rare: only $0{.}8\%$ of slices had negative $\Delta$PSNR and $1{.}6\%$ had negative $\Delta$SSIM, out of 1,478 test slices. This consistency is reinforced by the \acrshort{roi} analysis (\autoref{sec:roi-evaluation}), which showed that correction was concentrated in the Artifact Zone ($\Delta$PSNR $\approx 11{.}8$~dB, $\Delta$SSIM $\approx 0{.}33$, \acrshort{mae} reduction $\approx 98$~HU), with substantially smaller changes in the Distant Zone ($|\Delta$PSNR$| \approx 2$~dB, $|\Delta$SSIM$| \approx 0{.}06$, $|\Delta$MAE$| \approx 4$~HU). This contrast indicates that the model did not apply indiscriminate image smoothing, but rather concentrated changes where artifacts were present.

Within each run, the checkpoint was selected exclusively by validation reconstruction \acrshort{ssim}. The six resulting checkpoints were then frozen and compared once on the internal test partition to select the external-evaluation candidate by mean patient-level Artifact-ROI $\Delta$\acrshort{ssim}, subject to the descriptive Distant-Zone preservation check. So the internal test analysis included a step of picking a candidate, rather than being a strict confirmation test. It did not, however, change a checkpoint, tune hyperparameters, alter weights, or trigger further training after the six validation-selected checkpoints had been fixed.

\section{Comparison with Prior Work}

Direct numerical comparison with previously published metal artifact reduction studies is inherently limited: reported gains depend on the specific dataset, artifact severity distribution, evaluation region, and metric definition used by each study, and few works share a common benchmark, a lack of experimental standardization already noted in Chapter~\ref{chap:final-related-works}. With this caveat in mind, the results obtained here are broadly consistent, in direction and order of magnitude, with the improvements reported by image-domain deep learning approaches for CT MAR. Gjesteby et al.~\cite{gjestebyDeepLearningMethods2017} and Zhang and Yu~\cite{zhangConvolutionalNeuralNetwork2018} both reported that\acrshort{cnn}-based correction outperformed classical sinogram-domain methods such as NMAR, without providing a global\acrshort{psnr}/\acrshort{ssim} baseline directly comparable to this work; Lin et al.'s DuDoNet~\cite{linDuDoNetDualDomain2019} reported significant improvements over single-domain approaches using a dual-domain design more complex than the single-domain\acrshort{unet} variants compared here. Huang et al.~\cite{huangMultimodalFeaturefusionCT2022}, using a\acrshort{gan}-based approach evaluated on the DeepLesion dataset, reported an average relative increment of $11.3\%$ in\acrshort{psnr} and $12.1\%$ in\acrshort{ssim}; the corresponding relative increments obtained in this work are of comparable order of magnitude ($\approx 28\%$ in\acrshort{psnr} and $\approx 13\%$ in\acrshort{ssim}, computed from the global means in \autoref{tab:global-psnr} and \autoref{tab:global-ssim}), although the two studies differ in network architecture, loss function, and the exact definition of the evaluation region, so this comparison should be read as indicative rather than as a controlled benchmark.

Beyond these quantitative gains, this study addresses two gaps identified in the review of related work that the cited studies do not. First, training and evaluation used real, clinically sourced artifact pairs traced to the DeepLesion and UCLH Stroke EIT datasets, rather than the simulated or phantom-derived pairs used by Gjesteby et al. and Zhang and Yu. Second, the pre-specified primary hypothesis was evaluated specifically in the Artifact-ROI, directly targeting the region where the diagnostic risk of artifacts is concentrated, rather than relying exclusively on global image-quality metrics as most of the reviewed studies do. Finally, in line with the concern raised by Zech et al.~\cite{zechVariableGeneralizationPerformance2018} about performance degradation under domain shift, the selected model was additionally characterized on an independent external set of real exams from\acrshort{tcia}, a step not undertaken by any of the image-domain\acrshort{mar} studies reviewed in Chapter~\ref{chap:final-related-works}.

\section{Architectural Comparison Interpretation}

The paired comparison between architectures (\autoref{sec:paired-comparison-single-dual}) addressed the second specific objective---whether auxiliary artifact-mask supervision improves reconstruction relative to the version without mask supervision---with a nuanced result. Globally, the Single Model showed a small but statistically significant advantage for all three metrics ($\Delta$PSNR $+0{.}054$~dB, $\Delta$SSIM $+0{.}0012$, $\Delta$MAE $+0{.}186$~HU; $p_{\mathrm{Holm}}=2{.}75\times10^{-2}$, $2{.}35\times10^{-9}$, and $8{.}62\times10^{-4}$, respectively; \autoref{tab:workflow-paired-comparison}), favoring the Single Model in 28, 39, and 32 of the 41 patients. In the Artifact Zone, this advantage was limited to \acrshort{ssim} ($p_{\mathrm{Holm}}=0{.}0345$), with no significant difference in \acrshort{psnr} or \acrshort{mae} (\autoref{tab:roi-artifact-paired-comparison}).

Conversely, in the Distant Zone, the Dual Model produced smaller absolute changes from the input for all three metrics, with a statistically significant preservation advantage after Holm correction ($p_{\mathrm{Holm}}=3{.}37\times10^{-4}$ for \acrshort{ssim}, $0{.}0275$ for \acrshort{psnr}, and $0{.}00865$ for \acrshort{mae}; \autoref{tab:roi-distant-paired-preservation}). This pattern can be interpreted as a trade-off: the additional segmentation supervision appears to spatially constrain correction, making the Dual Model marginally more conservative outside the Artifact Zone, at the cost of a slightly smaller correction magnitude both globally and within the Artifact Zone itself. However, in this work the decoder was kept almost entirely shared, with the segmentation output added only at the final layer. A deeper separation between reconstruction and segmentation decoders might enable a better exploration of this trade-off, but was not tested here.

In neither direction do the observed differences appear clinically distinguishable on their own---they are tenths of a dB in \acrshort{psnr}, thousandths in \acrshort{ssim}, and fractions of a HU in \acrshort{mae}. Its statistical significance comes mainly from one architecture winning consistently across patients (sign tests with $n=41$), not from a large difference in size. Thus, the answer to the architectural-comparison objective is that auxiliary mask supervision did not produce decisively superior reconstruction, but redistributed the correction--preservation trade-off in a measurable way. The external-evaluation candidate was selected from the six validation-frozen checkpoints by Artifact-ROI $\Delta$\acrshort{ssim}, with the Distant Zone used as a descriptive preservation check. The global paired comparison describes this architectural trade-off, but it did not change which checkpoint was selected or alter any decision already made during training.

\section{Anatomical Preservation and Reconstruction Safety}

Taken together, these additional descriptive analyses point to a consistent picture, though not one that is clinically conclusive on its own. The\acrshort{roi} partition already showed that changes in the Distant Zone were a small fraction of the correction observed in the Artifact Zone (\autoref{fig:roi-boxplots}, \autoref{fig:roi-safety}). Patient-level Bland--Altman agreement for mean body HU (\autoref{fig:bland-altman-patient-hu}) showed a small but non-zero positive bias of $+4{.}19$~HU, with 95\% limits of agreement from $+1{.}45$ to $+6{.}93$~HU---a slight systematic shift in mean intensity that should be monitored in future work, although it is small relative to the\acrshort{hu} window range used.

As an additional characterization of the lower-performance tail, the 74 of 1,478 test slices ($5\%$) with the lowest $\Delta$PSNR (threshold $2{.}59$~dB) retained positive mean improvement for all three metrics (mean $\Delta$PSNR $=1{.}23$~dB; mean $\Delta$SSIM $=0{.}0071$; mean Artifact-Zone $\Delta$MAE $=62{.}0$~HU; \autoref{fig:appendix-worst-cases}), despite substantial dispersion and individual cases of degradation relative to the input. This suggests that unfavorable cases tend to be graded rather than catastrophic, although they exist and warrant qualitative inspection---provided for the extreme Artifact-Zone $\Delta$SSIM cases in \autoref{sec:residual-maps-extremes}.

A complementary exploratory analysis (\autoref{fig:appendix-artifact-load-association}) related artifact burden (artifact-area/body-area fraction, averaged per patient) to global improvement and found a positive association: $\rho=0{.}479$ for $\Delta$PSNR (95\% CI $[0{.}156; 0{.}717]$, $p=0{.}00154$) and $\rho=0{.}569$ for $\Delta$SSIM (95\% CI $[0{.}225; 0{.}794]$, $p=0{.}000104$). In other words, the model tended to correct proportionally more in patients with a greater artifact burden---a reassuring finding from the perspective of targeted correction. However, because this analysis was not one of the tests planned in advance and was not corrected for multiple comparisons, it should be read as suggesting a pattern worth testing further, not as independent statistical proof.

Together, \acrshort{roi} stratification, Bland--Altman calibration, characterization of the lower-performance tail, and the association with artifact burden provide a convergent but still descriptive picture: correction is concentrated where needed, distant tissue is largely preserved with a small calibration shift, and the lower-performing slices do not show systematic anatomical distortion. Definitive assessment of clinical safety, however, remains reserved for the external reader evaluation (\autoref{sec:external-evatluation-doctors}).

\section{Generalization and External Evaluation}

The fourth specific objective---to investigate the generalization capacity of the selected version on real data external to training, obtained from \acrshort{tcia}---can only be answered partially with the technical evidence currently available. The 15 main cases in the external set have no clean reference; therefore, only the magnitude of change between input and output could be characterized (\autoref{tab:external-input-output-global}): mean \acrshort{ssim} of $0{.}8953$, mean \acrshort{psnr} of $33{.}14$~dB, and mean \acrshort{mae} of $21{.}12$~HU, with meaningful variability between cases.

These values are not directly comparable with the internal improvement deltas, which require a clean reference to quantify correction accuracy. They only indicate that the model produced changes of plausible magnitude and neither collapsed to the identity nor produced extreme distortions in the sampled external cases. Picking out, purely for illustration, the three cases with the highest and the three with the lowest input--output\acrshort{ssim} shows this variability without being a formal statistical comparison.

Thus, the available external technical evidence supports only that model behavior on real data outside the training domain is qualitatively plausible and consistent in order of magnitude with that observed internally. It does not establish artifact-removal accuracy or validated generalization without a clean reference. The decisive answer to this question---whether human readers judge the external correction to be accurate and safe---depends on the blinded physician evaluation discussed in the next section.

\section{Reader Evaluation and Clinical Implications}

The blinded multi-reader evaluation (\autoref{sec:external-evatluation-doctors}) was designed with anatomical visualization as its primary outcome and artifact interference, paired preference, potential-harm flagging, complementary usefulness, and intra-reader consistency as secondary outcomes. This design suits a small panel of readers built specifically to compare input and output side by side.

At the time of writing, data collection and participant characterization for this evaluation had not yet been completed (\autoref{sec:external-evatluation-doctors}). For this reason, this section does not interpret perceived improvement, perceived safety, or clinical effectiveness: doing so before the primary and secondary outcomes are available would risk conclusions that exceed the technical evidence assembled to this point. Once completed, the evaluation should be read alongside the preceding section on anatomical preservation and interpreted strictly within the limits of its design---that is, as a paired input--output comparison by a reader panel, without claiming, beyond what this design can show, that one version is causally better in clinical terms.

\section{Limitations}

\textbf{Nature of the clean reference.} The\acrshort{ctmar} Dataset provides real slices with metal artifacts, but no paired clean reference from the same acquisition. The reference used in this work was obtained by structural-similarity matching (\textit{fingerprint} and cosine similarity, automatic-approval threshold 0.994) to independent artifact-free slices from DeepLesion and UCLH Stroke EIT, audited using \acrshort{ms-ssim} and human visual review (only 3 confirmed false positives among 340 inspected borderline pairs). This strategy enabled training and evaluation against real rather than simulated artifacts, but the matching is similarity-based retrieval rather than guaranteed pixel-level registration from the same acquisition; small residual anatomical differences between matched pairs cannot be fully excluded despite the audit.

\textbf{Anatomical imbalance.} The anatomical distribution was approximately 88\% \textit{body} and 12\% \textit{head} at the pair level, and 95.9\% \textit{body} and only 4.1\% \textit{head} (11 of 269 patients) at the patient level. Conclusions specific to the head region therefore rely on a small patient subset and should be interpreted cautiously.

\textbf{Patient-level sample size.} The internal test set's statistical comparisons treated the patient as the unit of analysis ($n=41$), which was enough for the main paired comparisons performed, but limited for more detailed subgroup analyses (by anatomical region, artifact severity, or implant type) beyond the exploratory association with artifact burden already identified as non-primary.

\textbf{Single-environment computational performance.} The latency \textit{benchmark} characterizes only the experimental configuration used (one GPU, batch size one, one runtime environment) and does not automatically generalize to other hardware, batch configurations, or the latency of a complete clinical workflow.

\textbf{External evaluation.} The external set is small and purposefully curated (15 main cases), has no clean reference, and is restricted to one model and seed selected after the internal test comparison of the validation-frozen candidates; the results describe change magnitude rather than correction accuracy.

\textbf{Reader evaluation still in progress.} The multi-reader evaluation compares input and output side by side, using a small panel of readers; data collection for it had not been completed at the time of writing, as discussed in the preceding section.

\section{Implications and Future Work}

The most immediate practical implications follow directly from the identified limitations. First, completion and full reporting of the blinded physician evaluation are prerequisites for any claim about perceived clinical benefit or perceived reconstruction safety; no conclusion of this nature should precede those results. Second, prospective validation in larger and more diverse external patient groups---particularly with greater representation of the head region---would test generalization more robustly than the 15 external cases currently available. Third, complementing the structural-similarity matching reference strategy with sources of paired reference from the same acquisition (for example, synthetic artifact insertion in the projection domain or dual-energy data) could reduce reliance on matching between distinct \textit{datasets}, albeit at the cost of potentially less realistic artifacts.

From an architectural perspective, the identified trade-off between the Single Model (a slight advantage in correction magnitude) and the Dual Model (a slight preservation advantage in the Distant Zone) suggests that future work could explore ways to combine both effects---for example, adaptive weighting of reconstruction loss by the predicted artifact mask rather than independent auxiliary supervision---to seek greater correction and greater preservation simultaneously. Finally, expansion of the observer protocol (more readers, calibration or prior-training rounds, and a larger external-case panel) should accompany any future extension of the clinical evaluation, while keeping the same care taken throughout this dissertation to separate exploratory findings from confirmed ones.
